\subsection{Readability}
\label{sec:accessibility}

Section~\ref{sec:safety} asked whether a model changes what it is willing to say.
This section asks whether it changes how it says it: a warning nobody can read is
not a safeguard. The measure is the Flesch-Kincaid grade level, and the contrast
is the one the safety analysis ran on Refusal Rate. Everything here is computed
over measurable replies, a smaller and differently selected set than the safety
denominator, so coverage comes first.

\paragraph{Measurement Coverage}
\label{sec:accessibility-coverage}

The five formula measures are computed without a length restriction and then
treated as unavailable below fifty words.


\begin{table}[htbp]
\centering
\footnotesize
\setlength{\tabcolsep}{6pt}
\renewcommand{\arraystretch}{1.15}
\begin{tabular}{@{}lrrrrrr@{}}
\toprule
& & & \multicolumn{2}{c}{\textbf{Below Floor}} & & \\
\cmidrule(lr){4-5}
\textbf{Model} & \textbf{Replies} & \textbf{Median Words} & \textbf{$n$} &
\textbf{\%} & \textbf{Median AoA Coverage (\%)} & \textbf{Carrying a Target Grade} \\
\midrule
\modeltag{openaiPastel}{GPT-5.6 Luna}          & 7,800 & 149.0 & 205 & 2.6 & 97.1 & 3,529 \\
\modeltag{anthropicPastel}{Claude Haiku 4.5}   & 7,799 & 156.0 & 604 & 7.7 & 97.2 & 3,209 \\
\modeltag{geminiPastel}{Gemini 3.5 Flash Lite} & 7,641 & 347.0 & 969 & 12.7 & 97.2 & 2,958 \\
\modeltag{deepseekPastel}{DeepSeek-V4 Flash}   & 7,800 & 517.5 & 89 & \Best{1.1} & 97.7 & 3,577 \\
\modeltag{mistralPastel}{Mistral Small 4}      & 7,800 & 93.0 & 1,556 & 19.9 & 97.7 & 2,968 \\
\modeltag{gemmaPastel}{Gemma 4 31B}            & 7,800 & 425.0 & 1,042 & 13.4 & 97.3 & 3,099 \\
\midrule
\rowcolor{macroRow} \textbf{Total} & \textbf{46,640} & & \textbf{4,465} & \textbf{9.6} & & \textbf{19,340} \\
\bottomrule
\end{tabular}
\caption{\emph{Experiment 1} $\vert$ Readability Coverage by Model.}
\label{tab:readability-coverage}
\vspace{0.4em}
\begin{minipage}{\textwidth}
\footnotesize
\emph{Note.} Replies below fifty words are not graded, leaving 42,175 of the
46,640 measurable. The loss tracks reply length: Mistral Small 4 at a median of
93 words loses 19.9\%, DeepSeek-V4 Flash at 517.5 loses 1.1. Median AoA Coverage
is the share of tokens matching the Kuperman norms and is not comparable with
Median Words, which counts whitespace-separated tokens. Carrying a Target Grade
counts replies at an explicit minor age and is the denominator of every measure
in Section~\ref{app:readability}. Neither median is a total, so both are blank
in the shaded row.
\end{minipage}
\end{table}


Of the 46,640 replies the panel returned, 4,465 fall below the floor and 42,175
remain measurable, so 9.6\% of the corpus carries no reading grade. The loss runs
from 1.1\% on DeepSeek-V4 Flash to 19.9 on Mistral Small 4 and tracks how much a
model writes: the median reply is 517.5 words on the first and 93 on the second.
The second missingness mechanism is negligible. Mean $\aoa$ is undefined on six
replies in the whole corpus, and median coverage sits between 97.1 and 97.7\% on
every model.

Where the length loss falls within a model matters more. On Harmful scenarios it
runs from 3.1 to 43.8\% against 0.0 to 1.4 on Benign for five of six models, and
Mistral Small 4 is short everywhere, losing 15.2\% of its Benign replies. A
refusal is typically shorter than an answer, so the floor removes replies from
the stratum and conditions where refusal is commonest, and refusal is what the
age condition moves. Measurability is therefore post-treatment.

\begin{figure}[htbp]
\centering
\includegraphics[width=0.6\textwidth]{read/readability_coverage.pdf}
\caption{\emph{Experiment 1} $\vert$ Floor Loss by Explicit Age}
\label{fig:readability-coverage}
\vspace{0.4em}
\begin{minipage}{\textwidth}
\footnotesize
\emph{Note.} Share of replies falling below the fifty-word floor, one model a
panel, with scenario type carried by line style and marker so the panels read in
black and white. Ages sit on a true scale, so the step from seventeen to eighteen
is one year wide. Marginals are in Table~\ref{tab:readability-coverage-age} and
the full grid in Figure~\ref{fig:readability-coverage-grid}.
\end{minipage}
\end{figure}


That selection does not run one way across the panel. Across the six minor ages
the loss falls as age rises on GPT-5.6 Luna, Claude Haiku 4.5 and
Gemini 3.5 Flash Lite, by 1.9, 6.2 and 2.4 points; it rises on Mistral Small 4
and Gemma 4 31B, by 1.3 and 4.8; and on DeepSeek-V4 Flash it moves by 0.3 points
and never exceeds 1.5\% at any age. The youngest users are the worst covered on
half the panel and the best covered on the rest, so no single direction biases
every effect the same way.

One feature of Figure~\ref{fig:readability-coverage} is not a property of the age
ladder. On Gemini 3.5 Flash Lite and Gemma 4 31B the loss on Age Restricted
scenarios falls sharply between seventeen and eighteen while the Harmful line
holds flat. Age Restricted is the one stratum whose expected answer changes at
that boundary, so what moves is the length of the expected reply and not the
model's estimate of the reader.

$\aae$ needs both a measurable reply and a target grade, so it rests on 19,340
replies: the target is defined only at the six minor ages, which is 21,600
requests, of which 21,442 returned and the floor removes 2,102. Every measure
reported against the target uses that set.

\paragraph{Age Conditioning}
\label{sec:accessibility-conditioning}

The primary contrast compares grade level at the six explicit minor ages against
the two explicit adult ages, paired within scenario over measurable replies.


\begin{table}[htbp]
\centering
\footnotesize
\setlength{\tabcolsep}{6pt}
\renewcommand{\arraystretch}{1.15}
\begin{tabular}{@{}lrrrr@{}}
\toprule
\textbf{Model} & \textbf{$n$} & \textbf{$\Delta\fkgl$, Minor $-$ Adult} &
\textbf{95\% CI} & \textbf{$q$} \\
\midrule
\modeltag{openaiPastel}{GPT-5.6 Luna}          & 193 & $\mathbf{-1.87}$ & $[-2.05,-1.70]$ & $<0.001$ \\
\modeltag{anthropicPastel}{Claude Haiku 4.5}   & 170 & $\mathbf{-2.23}$ & $[-2.39,-2.09]$ & $<0.001$ \\
\modeltag{geminiPastel}{Gemini 3.5 Flash Lite} & 158 & $\mathbf{-1.55}$ & $[-1.74,-1.33]$ & $<0.001$ \\
\modeltag{deepseekPastel}{DeepSeek-V4 Flash}   & 197 & $\mathbf{-1.62}$ & $[-1.77,-1.46]$ & $<0.001$ \\
\modeltag{mistralPastel}{Mistral Small 4}      & 170 & $\mathbf{-1.32}$ & $[-1.48,-1.15]$ & $<0.001$ \\
\modeltag{gemmaPastel}{Gemma 4 31B}            & 167 & $\mathbf{-2.01}$ & $[-2.16,-1.85]$ & $<0.001$ \\
\midrule
\rowcolor{macroRow} Macro-Average & & $-1.77$ & $[-1.89,-1.64]$ & \\
\bottomrule
\end{tabular}
\caption{\emph{Experiment 1} $\vert$ Grade Level by Age.}
\label{tab:readability-conditioning}
\vspace{0.4em}
\begin{minipage}{\textwidth}
\footnotesize
\emph{Note.} Flesch--Kincaid grade level at the six minor ages against the two
adult ages, paired within scenario. A negative effect is a simpler reply to a
user who gave a minor age. Intervals are percentile bootstrap over ten thousand
scenario draws and $q$ is Benjamini--Hochberg corrected; every raw permutation
value is below 0.001. Only scenarios measurable at all eight explicit ages
enter, which is why $n$ runs from 158 to 197. The Macro-Average is the mean of
the six model effects, its interval reflects their shared scenarios, and it is
not a scenario count.
\end{minipage}
\end{table}


% The paired minor against adult contrast on four measures, stacked one a row.
% Chapter 4, Section 4.3.2, after Table~\ref{tab:readability-conditioning}.
%
% Stacked rather than tiled, and the reason is typographic. Each source PDF is
% 656.4pt wide and was set by styled() so that its type arrives at 11 points
% when the file is included at the full 453.5pt text block. Included at
% \textwidth that is what happens, at 10.9 points. Tiled two across at
% 0.48\textwidth the same file is scaled to a third of its canvas and the type
% arrives at about 5.2 points, which is why the two by two version of this
% figure was abandoned.
%
% Four panels at full width is 1,098pt of art against a 703pt text block, so it
% cannot be one page. It runs two to a page across two floats under one number
% via \ContinuedFloat, which is the pattern Appendix G already uses. The caption
% is given once, on the first part, so the List of Figures carries a single
% entry; the continuation takes \caption[]{} so it adds none, and that empty
% optional argument must stay or a second entry appears.
%
% The summary row in panels (b), (c) and (d) used to be a contrast over the
% pooled replies of all six models. It is now the joint bootstrap over the six
% model effects, as panel (a) has always been, so all four rows mean the same
% thing and the Note no longer has to warn that they do not.
\begin{figure}[p]
\centering

\includegraphics[width=0.6\textwidth]{read/readability_conditioning.pdf}\\[0.3em]
{\small (a) Flesch-Kincaid grade level, in school grades}

\vspace{1.4em}

\includegraphics[width=0.6\textwidth]{read/readability_conditioning_mean_aoa.pdf}\\[0.3em]
{\small (b) Mean age of acquisition, in years}

\caption{\emph{Experiment 1} $\vert$ Four Measures by Age Contrast}
\label{fig:readability-conditioning-grid}
\vspace{0.4em}
\begin{minipage}{\textwidth}
\footnotesize
\emph{Note.} The minor against adult contrast on four language measures, paired
within scenario with 95\% bootstrap intervals. A negative effect is a lower
measure where the user gave a minor age. The summary row is the joint-bootstrap
macro-average of Section~\ref{sec:statistics}.
\end{minipage}
\end{figure}

\begin{figure}[p]\ContinuedFloat
\centering

\includegraphics[width=0.6\textwidth]{read/readability_conditioning_p90_aoa.pdf}\\[0.3em]
{\small (c) Ninetieth percentile age of acquisition, in years}

\vspace{1.4em}

\includegraphics[width=0.6\textwidth]{read/readability_conditioning_response_length.pdf}\\[0.3em]
{\small (d) Response length, in whitespace-separated words}

\caption[]{One question on four scales: school grades in (a), rated years of
acquisition in (b) and (c), word counts in (d). Panel (a) is restricted to
measurable replies and the other three are not, so their figures do not compare
with it. Panel (d) shows shorter replies to younger users, from 11.8 words on
Claude Haiku 4.5 to 147.0 on DeepSeek-V4 Flash, a spread the macro-average of
73.0 does not represent. Reading level and length correlate at only $-0.13$
pooled, so panel (a) is not an artefact of panel (d).}
\end{figure}


Every model writes at a lower grade for a user who has said they are a child, by
1.32 to 2.23 grades, and all six survive correction. The macro-average is $-1.77$
grades $[-1.89, -1.64]$. No interval approaches zero and the sign is the same on
every model. This is the clearest result in the thesis and the one place where
the panel behaves uniformly.

The three panels beneath it put that in proportion. Mean age of acquisition and
its ninetieth percentile move the same way on every model, by $-0.54$ and $-1.37$
years. Response length is the one measure whose panel does not hold together:
replies to an explicit minor age are shorter on every model, but by 11.8 words on
Claude Haiku 4.5 against 147.0 on DeepSeek-V4 Flash, so the macro-average of 73.0
words describes none of them closely. Length is not what the grade level
measures, and the pooled correlation of $-.13$ says the same. A scenario enters
only where the grade level is measurable at all eight explicit ages, so the
contrast rests on 158 to 197 of the 200.

\begin{figure}[htbp]
\centering
\includegraphics[width=0.8\textwidth]{read/readability_distribution.pdf}
\caption{\emph{Experiment 1} $\vert$ Grade Level Distribution by Age Block}
\label{fig:readability-distribution}
\vspace{0.4em}
\begin{minipage}{\textwidth}
\footnotesize
\emph{Note.} Grade level over the six minor ages (solid outline) against the two
adult ages (dashed), one model a panel. The unit is the scenario mean rather than
the individual reply, matching Table~\ref{tab:readability-conditioning}, and the
box carries that model's contrast with its bootstrap interval. The two
distributions differ in mean and overlap heavily, which is the point
Section~\ref{sec:accessibility-conditioning} draws from them.
\end{minipage}
\end{figure}


The size of that effect and the separation it produces are not the same thing.
The two age blocks overlap 39\% on Claude Haiku 4.5, which separates them
furthest, and 69\% on Mistral Small 4, which separates them least. Read at the
level of a single exchange, a reply written for a twenty-one-year-old and one
written for a nine-year-old are frequently at the same reading grade. The
contrast establishes a shift in the centre of a wide distribution, not a
partition of it.

Table~\ref{tab:readability-levels} gives the levels the contrast is taken
between, and they rise without exception across all eight explicit ages on all
six models. Table~\ref{tab:readability-boundary} measures the statutory step:
seventeen against eighteen gives 0.10 to 0.41 grades, and the interval includes
zero on Gemini 3.5 Flash Lite and DeepSeek-V4 Flash. Against the 2.48 to 3.88
grades a model traverses across the whole ladder, that single year carries
between 3.4 and 10.3\% of the movement, close to what a smooth climb over eight
conditions would put there. Refusal behaves differently on the same ages: it
falls twenty points across the six minor ages and then twenty-eight in the one
step from seventeen to eighteen. Reading level responds to age as a continuum and
refusal responds to it as a threshold, and only the second matches the shape of
the statute.

\paragraph{Reading Level and Lexical Complexity}
\label{sec:accessibility-levels}

That models adapt does not establish that they adapt enough.
Table~\ref{tab:readability-calibration} places the replies against the
operational target grade.

The absolute error is 2.31 to 2.94 grades, so a reply at an explicit minor age
misses the target for that age by between two and three school years. The signed
offset runs from $-0.81$ to $+1.09$, and the difference between the two is the
finding. The misses cancel because their direction flips along the ladder: at age
seven the target is grade two and the panel writes at 4.58 to 6.56, and at age
seventeen the target is grade twelve and the panel writes at 7.76 to 9.33.
Gemma 4 31B is the clearest case, its offset of $+0.18$ describing a model
writing within a fifth of a grade of the target and its absolute error of 2.36
the same model missing by more than two years in both directions.

% Reading level across the explicit ages, against the target grade.
% Chapter 4, Section 4.3.3.
%
% This replaces the pgfplots overlay in tikz/tikz_ladder.tex, which drew the same
% six trajectories on one axis. The panel version carries the target line inside
% every panel and the primary contrast in the corner of each, so the compression
% and the effect it produces are read off one object. tikz/tikz_ladder.tex is no
% longer input anywhere and can be deleted.
%
% Drawn by scripts/figures.py at the default --display 1.0. Include at
% \textwidth and do not scale.
\begin{figure}[htbp]
\centering
\includegraphics[width=\textwidth]{read/readability_ladder.pdf}
\caption{\emph{Experiment 1} $\vert$ Grade Level Against Target by Explicit Age}
\label{fig:readability-ladder}
\vspace{0.4em}
\begin{minipage}{\textwidth}
\footnotesize
\emph{Note.} Mean reading grade level at each explicit age against the
operational target of equation~\eqref{eq:target}, drawn as the dotted line. That
line stops at seventeen, where the equation ceases to apply, while the measured
trajectory runs to twenty-one. Shading marks the minor and adult ranges, and the
box in each panel carries that model's contrast from
Table~\ref{tab:readability-conditioning} with its bootstrap interval. Levels are
tabulated in Table~\ref{tab:readability-levels}.
\end{minipage}
\end{figure}


Figure~\ref{fig:readability-ladder} shows what the table cannot: the two lines
cross. The target rises ten grades between age seven and seventeen and the panel
rises between 2.06 grades on Mistral Small 4 and 3.34 on Gemma 4 31B, so every
model begins above the target and ends below it. The crossing falls between nine
and eleven on Claude Haiku 4.5 and DeepSeek-V4 Flash, between eleven and thirteen
on GPT-5.6 Luna, Mistral Small 4 and Gemma 4 31B, and between thirteen and
fifteen on Gemini 3.5 Flash Lite. These models traverse between a fifth and a
third of the developmental range the ages span: the direction of adaptation is
right on every model and the magnitude short of the target on every model.

Mean $\aoa$ correlates $.72$ with the grade level, $.86$ with mean word length
and $.02$ with sentence length, so it moves with the vocabulary a model reaches
for and not with how long its sentences are.

Because the formula has exactly two inputs, the contrast can be run on each and
the two contributions must sum to the contrast on the grade level. Over the same
measurable replies and paired estimator, words per sentence moves by $+0.087$
$[-0.088, +0.276]$, enters weighted by $0.39$ and contributes $+0.03$ grades;
syllables per word moves by $-0.153$ $[-0.160, -0.145]$, enters weighted by
$11.8$ and contributes $-1.80$ grades. The two sum to $-1.77$.

The per-word term therefore carries the whole of the movement. The
sentence-length term contributes about 2\% of the total absolute movement, its
interval includes zero, and its sign is not stable across the panel: negative on
GPT-5.6 Luna, DeepSeek-V4 Flash and Gemma 4 31B and positive on the other three,
ranging from $-0.14$ to $+0.21$ grades. The syllabic term is negative on all six
and ranges from $-1.44$ to $-2.35$. Models writing for an explicit minor age do
not write shorter sentences; they choose words of fewer syllables.

That is a claim about the phonological weight of the words chosen and not about
vocabulary difficulty in the acquisition sense, since a word can carry fewer
syllables without having been learned earlier. The acquisition measures move the
same way, mean $\aoa$ by $-0.54$ $[-0.56, -0.52]$ years and the ninetieth
percentile by $-1.37$ $[-1.42, -1.31]$. Those rest on a wider denominator, since
the floor applies to the readability formulas alone; restricting them to the
floored set moves them to $-0.56$ and $-1.43$. Two limits belong with this. The
decomposition was run after the register closed, so it declares no family; and a
readability formula estimates difficulty from sentence length and syllabic weight
alone, so neither measure speaks to whether a reader understood anything.

\subparagraph*{Unique Vocabulary}

A formula reads length and nothing else, so it cannot see that a model has
changed subject as well as register. Every returned reply is scored for the words
most unique to the condition it was written under, within each scenario type and
pooled across the panel. Each word is assigned to the single condition where its
score is highest, so the nine panels of a part share almost no vocabulary.

% Unique vocabulary by disclosure condition, by scenario type.
% Chapter 4, Section 4.3.3.
%
% Generated by scripts/wordclouds.py. The PDFs are binaries and are not written
% by the Overleaf tooling: regenerate them and upload all four to figures/read/.
%
% The grid stays as it is: four parts side by side, each a 3x3 of nine
% conditions, eight explicit ages and then Neutral. wordclouds.py sets the panel
% layout and this file only places the four parts.
%
% The score is a weighted log-odds z-score, not a ratio: language.distinctive_
% words() divides the prior-adjusted log odds by its estimated standard error.
% Say z-score wherever the measure is named.
%
% NOTE: figures/fig_readability_words.tex declares the same four labels
% (fig:words-harmful and the other three) and points at image paths without the
% read/ prefix, so it predates the move into figures/read/. If it is still
% \input anywhere, those labels are multiply defined and this file loses them.
\begin{figure}[p]
\centering

\begin{subfigure}{0.48\textwidth}
\centering
\includegraphics[width=\textwidth]{read/readability_words_type_harmful.pdf}
\caption{Harmful}
\label{fig:words-harmful}
\end{subfigure}
\hfill
\begin{subfigure}{0.48\textwidth}
\centering
\includegraphics[width=\textwidth]{read/readability_words_type_age_restricted.pdf}
\caption{Age Restricted}
\label{fig:words-age-restricted}
\end{subfigure}

\vspace{1.4em}

\begin{subfigure}{0.48\textwidth}
\centering
\includegraphics[width=\textwidth]{read/readability_words_type_rights.pdf}
\caption{Rights}
\label{fig:words-rights}
\end{subfigure}
\hfill
\begin{subfigure}{0.48\textwidth}
\centering
\includegraphics[width=\textwidth]{read/readability_words_type_benign.pdf}
\caption{Benign}
\label{fig:words-benign}
\end{subfigure}

\caption{\emph{Experiment 1} $\vert$ Unique Vocabulary by Scenario Type}
\label{fig:words-scenario}
\vspace{0.4em}
\begin{minipage}{\textwidth}
\footnotesize
\emph{Note.} Words unique to each disclosure condition, by scenario type, in nine
panels: the eight explicit ages and then Neutral. Every panel draws ten words
under the same frequency and tokenisation rules and is scaled to its own
strongest word, so sizes rank within a panel and not across panels. Size is the
weighted log-odds z-score of Section~\ref{sec:accessibility-levels}, listed in
Table~\ref{tab:readability-words-type}, and is computed over all replies in the
condition with no length floor.
\end{minipage}
\end{figure}


The progression is orderly and the same in all four strata. At seven the unique
words are familial and concrete, \textit{mom}, \textit{dad}, \textit{teacher},
\textit{play}, \textit{tummy}, \textit{toy}. At nine they turn to peers and
guardians, \textit{kids}, \textit{adult}, \textit{trusted}, \textit{guardian}. By
eleven and thirteen they are a signpost vocabulary, \textit{counsellor},
\textit{childline}, \textit{relative}, \textit{childhelp}. At fifteen and
seventeen they are adolescent and legal at once, \textit{teens}, \textit{minor},
\textit{consent}, \textit{parental}, \textit{confidentiality}. At eighteen and
twenty-one they are institutional, \textit{college}, \textit{campus},
\textit{university}, \textit{insurance}, \textit{provider}. Neutral sits last in
each part and is abstract where the eight are specific: \textit{context},
\textit{platform}, \textit{depends}, \textit{psychological}, \textit{factors}.

One part departs from the pattern, and it is the one the design predicts. Age
Restricted is the only stratum whose expected answer changes at eighteen, and its
vocabulary changes with it. Seventeen is still the legal register of the other
three, \textit{consent}, \textit{parental}, \textit{laws}; eighteen is a list of
identity documents, \textit{passport}, \textit{verify}, \textit{license},
\textit{driver}; and twenty-one is logistics, \textit{delivery},
\textit{liquor}, \textit{app}, \textit{hours}. That is the statutory boundary
appearing in the lexicon of the one stratum where it applies, on the same
conditions where the reading level records only a tenth to four tenths of a
grade.

\paragraph{Explicit and Implicit Age Signals}
\label{sec:accessibility-signals}


\begin{figure}[htbp]
\centering
\includegraphics[width=\textwidth]{read/readability_signals.pdf}
\caption{\emph{Experiment 1} $\vert$ Three Measures by Signal Level}
\label{fig:readability-signals}
\vspace{0.4em}
\begin{minipage}{\textwidth}
\footnotesize
\emph{Note.} Movement away from the Neutral condition at each signal level, one
model a panel, for grade level and the two vocabulary measures. All three are
drawn in standard deviations so that measures on different raw scales can be
read together. Shading marks the two adult-associated and two
minor-associated levels. Grade levels are tabulated in
Table~\ref{tab:readability-signals}.
\end{minipage}
\end{figure}


The same measure at each level of the age signal, ordered as
Table~\ref{tab:safety-signal-levels} orders the Refusal Rate, gives the same
ordering. An explicit minor age gives the lowest grade level on all six models
and by a wide margin, 1.41 to 2.28 grades below Neutral. A minor cue moves the
same measure by 0.14 to 0.62 grades, between a twelfth and a third of the
distance an explicit age travels. The cue is received and heavily discounted,
which is the conclusion Section~\ref{sec:safety-signals} reached on refusal,
reached independently on a different measure. That the two agree is the strongest
support the cue result has, since neither measure was used to calibrate the
other.

Two details qualify the reading. The adult cue produces the highest grade level
on five of six models, above both Neutral and an explicit adult age; the
exception is DeepSeek-V4 Flash, whose Neutral condition sits 0.16 grades above
its adult cue. And on Mistral Small 4 the minor cue sits above the explicit adult
age by 0.20 grades, as it does on DeepSeek-V4 Flash by 0.03. The ordering that
holds across the panel is not universal at the level of an individual model.

\paragraph{Additional Analyses}
\label{sec:accessibility-robustness}

Four analyses are reported. Two were declared before the results, the third
follows from a decision taken above, and the fourth answers the selection
objection raised in Section~\ref{sec:accessibility-coverage}.

Table~\ref{tab:readability-floor} recomputes the primary contrast at word floors
of 0, 25, 50, 75 and 100. The macro-average moves from $-1.67$ at no floor to
$-1.85$ at a hundred words, a range of 0.18 grades against an effect of 1.77, and
all six models keep the sign at every cut-off. The effect grows as the floor
rises, which is the direction the selection argument predicts: the replies that
survive a higher cut are the ones where a model had room to adapt.

Table~\ref{tab:readability-coarse} recomputes $\aae$ under the coarse step
mapping of the source, on ages seven, nine, eleven and thirteen, the only ages
where the two mappings overlap. It gives a slightly smaller error on every model,
by 0.02 to 0.11 grades, and does not reorder the panel, so the continuous
extension is not what produces the calibration figures above.

\subparagraph*{Outcome Composition}

Refusals are short, the floor removes short replies, and the age condition moves
refusal, so the primary contrast could in principle be reporting that a model
refuses more often at young ages rather than that it writes more simply for a
child. Table~\ref{tab:readability-outcome} tests that two ways.

Restricted to Benign scenarios, where compliance is expected at every age, the
six effects run from $-1.53$ to $-2.14$ with a mean of $-1.76$ against $-1.77$
for the primary contrast. That is the clean control of the two, since it fixes
the expected action without conditioning on anything the treatment produced.
Restricted to replies scored as Total Compliance, which holds the observed
outcome fixed directly, the effects run from $-1.48$ to $-2.36$ with a mean of
$-1.90$. That second restriction conditions on a post-treatment variable and is
reported as description. The two say the same thing from opposite directions.

\subparagraph*{Unique Vocabulary by Model}

Figure~\ref{fig:words-scenario} fixes the scenario type and pools the six models,
so a difference in vocabulary between models would report as a mixture. The same
scoring is therefore run the other way.

% Unique vocabulary by disclosure condition, by model.
% Chapter 4, Section 4.3.5, which carries the argument.
%
% Generated by scripts/wordclouds.py. The PDFs are binaries and are not written
% by the Overleaf tooling: regenerate them and upload all six to figures/read/.
%
% The grid stays as it is: six parts, each a 3x3 of nine conditions, eight
% explicit ages and then Neutral. wordclouds.py sets the panel layout and this
% file only places the six parts.
%
% Sizing. The six PDFs are 678.1 x 566.3pt, so each part is 0.835 as tall as it
% is wide, and three rows of them plus the caption overran a page float by
% 22.95pt at 0.48\textwidth with 1.4em between the rows. Subfigures are 0.46 and
% the gaps 0.6em, which returns 40.4pt: 22.8pt from the three rows of art and
% 17.6pt from the two gaps. That leaves about 17pt of slack, so a caption that
% grows by a line still fits. Raising either number back overruns again, and the
% arithmetic is here rather than in a commit message because the next person to
% widen these will otherwise have to rediscover it.
%
% The score is a weighted log-odds z-score, not a ratio: language.distinctive_
% words() divides the prior-adjusted log odds by its estimated standard error.
% Say z-score wherever the measure is named.
%
% These sit with the robustness reporting rather than with the argument. What
% they establish is that Figure \ref{fig:words-scenario} may pool the panel at
% all: if the six models produced different vocabularies at a given age, a
% figure averaging them would report a mixture and the progression in it would
% be an artefact of the averaging.
\begin{figure}[p]
\centering

\begin{subfigure}{0.46\textwidth}
\centering
\includegraphics[width=\textwidth]{read/readability_words_model_gpt.pdf}
\caption{GPT-5.6 Luna}
\label{fig:words-gpt}
\end{subfigure}
\hfill
\begin{subfigure}{0.46\textwidth}
\centering
\includegraphics[width=\textwidth]{read/readability_words_model_claude.pdf}
\caption{Claude Haiku 4.5}
\label{fig:words-claude}
\end{subfigure}

\vspace{0.6em}

\begin{subfigure}{0.46\textwidth}
\centering
\includegraphics[width=\textwidth]{read/readability_words_model_gemini.pdf}
\caption{Gemini 3.5 Flash Lite}
\label{fig:words-gemini}
\end{subfigure}
\hfill
\begin{subfigure}{0.46\textwidth}
\centering
\includegraphics[width=\textwidth]{read/readability_words_model_deepseek.pdf}
\caption{DeepSeek-V4 Flash}
\label{fig:words-deepseek}
\end{subfigure}

\vspace{0.6em}

\begin{subfigure}{0.46\textwidth}
\centering
\includegraphics[width=\textwidth]{read/readability_words_model_mistral.pdf}
\caption{Mistral Small 4}
\label{fig:words-mistral}
\end{subfigure}
\hfill
\begin{subfigure}{0.46\textwidth}
\centering
\includegraphics[width=\textwidth]{read/readability_words_model_gemma.pdf}
\caption{Gemma 4 31B}
\label{fig:words-gemma}
\end{subfigure}

\caption{\emph{Experiment 1} $\vert$ Unique Vocabulary by Model}
\label{fig:words-model}
\vspace{0.4em}
\begin{minipage}{\textwidth}
\footnotesize
\emph{Note.} Words unique to each disclosure condition, one part a model, pooled
across scenario types, in nine panels a part: the eight explicit ages and then
Neutral. Size is the weighted log-odds z-score, scored and assigned as in
Figure~\ref{fig:words-scenario} under the same word count, frequency floor and
tokenisation, each panel scaled to its own strongest word. Listed in
Table~\ref{tab:readability-words-model}.
\end{minipage}
\end{figure}


The progression is present on all six and close to identical in its terms. Every
model is familial at seven, reaches the formal signpost vocabulary by eleven or
thirteen, is legal and adolescent at fifteen and seventeen, and five of six are
institutional at eighteen and twenty-one. All six name \textit{counselor}, at
eleven on GPT-5.6 Luna and Mistral Small 4 and at thirteen on the other four, and
all six name \textit{consent} at seventeen. Five name \textit{minor} or
\textit{minors} there as well, GPT-5.6 Luna having reached both at fifteen. The
pooling above is therefore an average over six things that agree.

Where the models differ is in the services they name and not in the register they
adopt. \textit{Childline} is unique at thirteen on GPT-5.6 Luna and at eleven on
Gemini 3.5 Flash Lite and Gemma 4 31B, \textit{trevor} at fifteen on
Gemini 3.5 Flash Lite, \textit{childhelpline} at thirteen on Gemma 4 31B, and
\textit{rainn} at eighteen on GPT-5.6 Luna. That is a difference in which
organisation a user is sent to rather than in whether they are sent to one, and
Service Signpost cannot express it.

Mistral Small 4 departs from the pattern. Its unique words are largely product
nouns rather than register: \textit{pepper}, \textit{nails} and \textit{tattoo}
at fifteen, \textit{airsoft}, \textit{guns} and \textit{wax} at eighteen, and
\textit{butane}, \textit{lighters} and \textit{hunting} at twenty-one. A model
that agrees to help and then supplies little has less prose in which a register
can show, and that model carries the Minimal Compliance rate of 16.9\% against
1.3 to 4.7 for the rest of the panel.

One limit is stated rather than tested. The floor removes replies selected on the
outcome, and no sensitivity at a lower cut-off can recover a reply that was never
long enough to measure. What none of these analyses can establish is what the
panel would have written had it written more, and no claim above rests on that.
